\documentclass[11pt]{article}
\usepackage[margin=1in]{geometry}
\usepackage{amsmath,amssymb}
\usepackage{booktabs}
\usepackage{hyperref}
\title{Empirical Design Notes: Sample Construction, Units, and Event
Exclusions\\
\large Rate-Based Reference Model (Direction F)}
\author{}
\date{}
\begin{document}
\maketitle

\section*{Scope}

This note documents the panel construction for the cross-sectional test of
$\phi^- = 1 - (\mathrm{ROE} - g)$: the units basis (Section~\ref{sec:units},
where an error in an earlier draft is corrected), what the regressor's
variance decomposition implies about what the test actually identifies
(Section~\ref{sec:decomp}), and the principled grounds for excluding
specific country-episodes (Section~\ref{sec:exclusions}). Every exclusion
is derived from a stated model assumption failing, not from an observation
being inconvenient for the predicted slope.

\section{Data}

The buffer $z$ (regulatory capital to risk-weighted assets, IMF FSI code
\texttt{FSKRC\_PT}) and the required return $\rho$ (approximated by return
on equity, \texttt{FSERE\_PT}) are drawn from the IMF Financial Soundness
Indicators dataset at quarterly frequency, Core Set, Deposit Takers.
Growth $g$ is drawn from the IMF World Economic Outlook database.

\subsection{Units: the regressor must be computed on a nominal basis}
\label{sec:units}

FSI return on equity is reported in nominal, domestic-currency terms: it
is net income (a nominal flow) divided by book equity (a nominal stock).
Required capital grows with the nominal balance sheet, since risk-weighted
assets are carried at nominal values. The regressor $\mathrm{ROE} - g$
is therefore only correctly specified when $g$ is \emph{nominal} GDP
growth in the same domestic currency --- not real GDP growth.

An earlier version of this analysis used WEO's \texttt{NGDP\_RPCH} (real,
constant-price growth) throughout. This was an error, not a stylistic
choice, and its effect is not small. Recomputing $g$ from WEO's
\texttt{NGDP} (current-price GDP level, domestic currency) as
$g_t = \mathrm{NGDP}_t/\mathrm{NGDP}_{t-1} - 1$ and averaging over each
country's sample window changes the regressor materially:

\begin{center}
\begin{tabular}{@{}lcc@{}}
\toprule
& real-basis (wrong) & nominal-basis (correct) \\
\midrule
mean of $\mathrm{ROE}-g$ & 10.34pp & 4.31pp \\
sd of $\mathrm{ROE}-g$    & 5.12pp  & 7.14pp \\
range                     & $-2.95$ to $24.00$ & $-34.42$ to $16.53$ \\
countries with negative rebuild capacity & 3 & 9 \\
\bottomrule
\end{tabular}
\end{center}

The starkest illustration is Argentina, whose regressor value moves from
$+24.00$ (the largest positive value in the real-basis panel) to
$-34.42$ (the largest negative value in the nominal-basis panel) --- a
58-point swing driven entirely by which growth series is used, for one
country. Turkey, Belarus, Angola and Ghana show swings of 15--26 points
for the same reason. The implied inflation proxy
($g_{\text{nominal}} - g_{\text{real}}$) exceeds 10 percentage points for
seven of seventy-five countries and has a panel median of 4.2pp, so this
is not confined to a handful of extreme cases; it is a first-order
correction to the regressor for a meaningful share of the sample.

\textbf{All results in this note and in any subsequent estimation use the
nominal basis.} The real-basis numbers should not be reported, including
as a robustness check, since they are not an alternative specification of
the same question --- they are the wrong quantity for the model as
written.

\section{What the regressor's variance decomposition implies}
\label{sec:decomp}

Across the panel (nominal basis, $n=75$), $\mathrm{ROE}$ has standard
deviation 4.95pp and $g$ has standard deviation 1.62pp (real) or
substantially more on a nominal basis in the high-inflation members; the
correlation between $\mathrm{ROE}$ and real $g$ is 0.06. Decomposing the
regressor's variance,
$\operatorname{Var}(\mathrm{ROE}-g) = \operatorname{Var}(\mathrm{ROE}) +
\operatorname{Var}(g) - 2\operatorname{Cov}(\mathrm{ROE},g)$, cross-country
variation in $\mathrm{ROE}$ accounts for the large majority of the
regressor's spread on the (mis-specified) real basis; on the nominal
basis, variation in $g$ itself is substantially larger, since nominal
growth ranges from near zero to over 50\% across the panel's
high-inflation members.

Two consequences follow.

\emph{First}, whichever basis is used, the paper's identifying variation
is not concentrated in near-zero-growth, low-inflation economies of the
kind (Japan, stagnant Western Europe) that motivate the theory
narratively. That mismatch between motivating case and estimating sample
must be stated explicitly rather than left for a reader to notice.

\emph{Second}, and more consequential on the nominal basis: the regressor
now spans very high nominal growth rates in chronically high-inflation
economies (Argentina, Turkey). Whether the linear model
$\phi^- = 1-(\mathrm{ROE}-g)$ is a sensible description of buffer dynamics
in a hyperinflationary or near-hyperinflationary economy is not
established by anything in this project and should not be assumed; these
observations may need their own treatment (Section~\ref{sec:exclusions})
rather than being pooled with the rest of the panel at face value.

\subsection{Negative rebuild capacity}

Nine countries (nominal basis) have $\mathrm{ROE} < g$, implying
$\phi^- > 1$: an explosive buffer under the model's linearisation. This is
economically a coherent prediction --- a banking system that cannot fund
its own required growth from earnings needs continuous external capital
injection to avoid a widening deficit --- but it is not estimable as an
ordinary autoregressive coefficient, and four of the nine (Argentina,
Turkey, Belarus, Tajikistan) are also flagged above for basis-sensitivity.

\paragraph{Decision, fixed before estimation (31 July 2026).} All nine are
treated as a \textbf{separate documented subsample}, reported alongside
the main estimation rather than pooled into it or dropped. Two of the
three options considered are rejected on stated grounds: an estimator
accommodating $\phi^- \geq 1$ is not adopted, because no such estimator
has been validated on known truth in this project --- standing practice
throughout this project is to validate any estimator against simulated
data of known truth before applying it to real data ---
and adopting one now, for exactly the
observations that would otherwise embarrass a linear fit, is the kind of
choice this document exists to prevent; outright exclusion under
Section~\ref{sec:exclusions} is not adopted either, because negative
rebuild capacity is not one of the four violations of B1--B2 listed
there --- it is a prediction of the model at these parameter values, not
a breakdown of its premises, and excluding a model's own predicted
regime because it is awkward to estimate is not a defensible exclusion
ground. A separate subsample is the only option that reports the
observation rather than discarding or mis-estimating it. The main-sample
result and the nine-country subsample result must both be reported;
neither may be presented alone.


\section{The deleveraging channel: B2 is incomplete, but the omission is
not growth-correlated}
\label{sec:deleveraging}

Assumption B2 states that a capital deficit can only be closed by earning
it back. That is incomplete: a deficit can also be closed by shrinking the
denominator --- shedding risk-weighted assets. Post-2008 European bank
capital ratios recovered substantially through balance-sheet contraction
rather than through profitability, so this is not a hypothetical channel.

The concern this raised was sharper than mere incompleteness. If
deleveraging is \emph{cheaper when growth is low} --- shedding assets costs
market share in a boom and much less when every competitor is shrinking
--- then the omitted channel would be correlated with $g$, and the
regression of $\phi^-$ on $(\mathrm{ROE}-g)$ would suffer omitted-variable
bias in a known direction: steeper than $-1$. That is the same direction as
the small-sample AR bias already documented, so the two would be
confounded and neither could be identified from the other.

\paragraph{Test.} Using IMF FSI risk-weighted-asset levels
(\texttt{FS\_ODX\_ARW\_FSKRC\_XDC}, available for 66 of the 77 panel
countries with clean overlap), quarterly RWA growth was computed and split
by regime, with the deficit indicator taken from the buffer ratio. The
statistic is the gap: surplus-quarter RWA growth minus deficit-quarter RWA
growth, so a positive gap means RWA grows more slowly (or contracts) when
the buffer is under pressure.

\paragraph{Result 1: the channel is real.} Sixty-four percent of countries
show a positive gap, with a cross-country mean of $1.19$ percentage points
per quarter. Banks do shrink risk-weighted assets when capital is short.
B2 as written omits something that happens.

\paragraph{Result 2: it is not growth-dependent, and the identification
concern does not survive.} The correlation between a country's average
nominal growth and its deleveraging gap is $+0.33$ --- the \emph{opposite}
sign to the one that would generate the feared bias. That weak positive
association is itself fragile: it falls to $+0.22$ once Eswatini and
Argentina are excluded (both already flagged in
Section~\ref{sec:exclusions}; Eswatini's gap of $19.7$ is more than double
any other country's, and Argentina's nominal growth is almost entirely
inflation), and the rank correlation --- which prevents two extreme
observations from driving the estimate --- is $0.19$ ($p = 0.13$) on the
full sample and $0.16$ ($p=0.20$) excluding those two. A median split on
growth is the clearest statement: low-growth countries average a gap of
$1.17$, high-growth countries $1.21$. Indistinguishable.

\paragraph{Consequence.} The deleveraging channel is a
\textbf{model-completeness} issue, not a \textbf{slope-identification}
threat. It belongs in the theory as an extension --- a deficit has a slow
cheap route (earn) and a fast costly one (deleverage), with the mix
depending on how binding the shortfall is --- but it does not need to be
resolved before the main cross-sectional estimation can be trusted, and the
overshoot to $-1.2/-1.3$ observed in simulation remains attributable to
small-sample AR bias rather than to this omission. This corrects a concern
raised earlier in the project's development, which anticipated the opposite
finding.

\paragraph{Coverage note.} Ten panel countries have no RWA level series at
all: Ecuador, Estonia, Croatia, Latvia, Malta, Netherlands, Palestine,
Slovakia, \textbf{the United States}, and Kosovo. This is a coverage gap on
the outcome variable, layered on the gaps already documented for the
regime variables, and the absence of the United States in particular should
be stated rather than passed over.

\paragraph{Eswatini.} The largest single outlier ($19.7$pp gap, more than
twice the next) was investigated rather than left unexplained. No
documented systemic banking crisis or bank resolution was found. What does
exist is a well-documented 2010--11 fiscal crisis --- a collapse in
Southern African Customs Union revenue --- at the very start of Eswatini's
FSI window, combined with an IMF assessment that Eswatini's banks carry
large exposures to the government such that a fiscal shock propagates
directly into the financial sector. This is a sovereign shock transmitted
to a very small banking system, not a banking crisis of the kind
Section~\ref{sec:exclusions} enumerates. It is therefore a candidate for
the size-based screen rather than the event-based one, and is flagged for
manual review rather than automatic exclusion.


\section{Estimating specification: regional fixed effects, with pooled OLS
reported alongside}
\label{sec:spec}

Comparing a small sub-Saharan banking system directly against a large
European one attributes to the model's mechanism what may simply be
structural difference between economies operating under different
pressures. The main specification therefore includes regional fixed
effects, identifying the slope from \emph{within}-region variation in
$(\mathrm{ROE}-g)$:
\[
  \hat\phi^-_i = \alpha_{r(i)} + \beta\,(\mathrm{ROE}_i - g_i)
  + \varepsilon_i,
  \qquad \text{model prediction } \beta = -1 .
\]
Pooled OLS without regional effects is reported alongside, so a reader can
see directly whether the result depends on cross-regional comparison.

\paragraph{Grouping.} Regions are assigned on \textbf{geographic proximity
alone}. Political affiliation, trade arrangements, currency blocs and
alliance structures are deliberately ignored; the grouping is meant to
capture the shared economic environment that comes with location, not
institutional membership. Eleven regions result:

\begin{center}
\begin{tabular}{@{}lrrr@{}}
\toprule
Region & $n$ & mean $(\mathrm{ROE}-g)$ & within-region sd \\
\midrule
Europe        & 19 & $3.02$  & $4.22$ \\
Sub-Saharan Africa & 16 & $6.69$  & $5.47$ \\
Latin America & 14 & $3.15$  & $11.23$ \\
Southeast Asia & 6 & $4.93$  & $3.02$ \\
Oceania       & 5  & $10.88$ & $4.82$ \\
West Asia     & 3  & $-2.25$ & $6.35$ \\
South Asia    & 3  & $0.99$  & $7.94$ \\
East Asia     & 3  & $6.65$  & $3.48$ \\
Middle East   & 2  & $7.31$  & $0.82$ \\
North America & 2  & $8.78$  & $5.37$ \\
Central Asia  & 2  & $-8.81$ & $4.87$ \\
\bottomrule
\end{tabular}
\end{center}

\paragraph{What fixed effects cost.} Decomposing the regressor's variance,
$24.9\%$ is between-region and $75.1\%$ within. Regional effects therefore
discard about a quarter of the identifying variation, leaving three
quarters --- ample for the slope test --- at a cost of eleven degrees of
freedom ($n=75$ falls to $63$ residual df). No region is a singleton, so
every observation contributes.

\paragraph{Three grouping decisions, stated rather than buried.} Georgia,
Armenia and Turkey are land-contiguous with both Europe and the Middle
East; rather than force a choice they are grouped as West Asia. Korea is
placed with Hong Kong and Macao as East Asia, not Southeast Asia. Bhutan,
Sri Lanka and the Maldives form South Asia; the region is not labelled
``India'' because India is not in the panel. Each of these would shift the
relevant regional means, and any of them is arguable; they are recorded here
so the choice is auditable rather than implicit.

\paragraph{An incidental benefit, and a cost.} Geographic grouping absorbs
some concerns that had been handled by separate rules. Latin America's
within-region standard deviation of $11.23$ --- far the largest, and driven
mostly by Argentina --- means regional effects partially absorb the
high-inflation problem of Section~\ref{sec:units} without a dedicated
exclusion. Central Asia separating cleanly from the Middle East
($-8.81$ against $+7.31$) is the grouping working as intended: an earlier
ad-hoc grouping had averaged those together and concealed the difference.

The cost is specific and should not be glossed. With between-region
variation removed, the specification can no longer test whether
structurally slower-growing \emph{regions} carry systematically higher
persistence. Given the growth distribution documented in
Section~\ref{sec:decomp} --- the panel contains almost no genuine
stagnation cases --- little is lost in practice, but the cross-regional
form of the stagnation claim becomes untestable by construction rather
than merely underpowered. That is why pooled OLS is reported alongside
rather than discarded.

\section{Economic grounds for excluding specific episodes}
\label{sec:exclusions}

Assumptions B1--B2 describe a going concern whose capital adjusts only
through retained earnings (deficit) or a payout decision (surplus). Four
situations are identifiable, in advance, as violations of this description
rather than as instances of it:

\begin{enumerate}
\item \textbf{External recapitalisation.} Where a shortfall is closed by a
  public injection, a bailout facility, or a state guarantee converted to
  equity, the deficit closes by fiat, not by earning it back. This is the
  single most consequential exclusion, because it is \emph{not neutral}: a
  recapitalised deficit closes fast without being earned, which pushes the
  observed $\hat\phi^-$ down. Removing these episodes therefore
  \emph{raises} estimated persistence and flatters the model's prediction
  at low growth --- the opposite of a conservative adjustment, and it must
  be reported as such.

\item \textbf{Payout restriction.} Where dividends and buybacks are
  prohibited by supervisory action --- the near-universal 2020--21
  restriction across advanced-economy banking systems being the largest
  instance in this sample window --- $\delta_p$ is set to zero externally,
  so $\phi^+ = 1$: surpluses become perfectly persistent, exactly the
  reverse of B2's assumption that a surplus is a policy choice.

\item \textbf{Resolution or nationalisation.} Where a bank is resolved,
  nationalised, or subject to an imposed balance-sheet contraction, the
  going-concern premise behind B1's growth identity does not hold.

\item \textbf{Definitional discontinuity.} The Basel~II to Basel~III
  transition changed what counts as regulatory capital, and adoption dates
  differ by jurisdiction, producing a level break in $z$ at a
  country-specific date unrelated to any economic event.

\end{enumerate}

\paragraph{A fifth, added after the units correction: chronic high
inflation.} Where a country's nominal growth rate is dominated by
inflation rather than real activity (Argentina and, over parts of the
sample window, Turkey and Belarus), the linear relationship
$\phi^- = 1-(\mathrm{ROE}-g)$ conflates a real economic mechanism with a
nominal accounting phenomenon that the model does not describe. These are
candidates for a separate high-inflation subsample rather than for the
main estimation, pending a decision on whether the model should be
respecified in real terms with a correspondingly real measure of
$\mathrm{ROE}$.

\paragraph{Protocol.} Event windows and the high-inflation flag are fixed
from external, pre-existing sources --- Laeven and Valencia's systemic
banking crisis database, in its current vintage extending through 2025
(IMF Working Paper WP/26/94, May 2026), for the first four grounds; a
simple threshold on average CPI inflation, fixed before estimation, for
the fifth --- \emph{before} any regression is run. Results are reported
with the flagged episodes and countries included and excluded, side by
side. A rule chosen after seeing its effect on the estimated slope is not
a rule.

\paragraph{What this is not.} The country-level amplitude screen used
earlier in sample construction (excluding, on the basis of return-on-equity
volatility alone, observations such as Ukraine during the war, the 2013
Cyprus bail-in, and the Greek banking crisis) conflates measurement noise
from a national ratio dominated by one or two institutions with genuine
systemic stress. The former is a data-quality problem and a candidate for
exclusion; the latter is exactly the kind of episode a model of
capital-buffer dynamics under stress ought to be tested against, not
screened away from. Where volatility is high for an identifiable systemic
reason, the preferred treatment is a robust estimator that limits its
leverage, or an explicit crisis subsample reported alongside the main
result, not deletion.

\section{The country universe}
\label{sec:universe}

Earlier versions of this panel drew on a fixed list of 77 areas. That list
excluded roughly half of the areas reporting the capital ratio, including
several with more usable history than areas it retained, and no rationale
for it is recorded anywhere in this project. Two candidate explanations
were tested directly against the source and both were rejected: it is not
the set of areas reporting both the capital ratio and return on equity
(those two indicators are reported by an identical set of areas), and it
does not correspond to any standard advanced-versus-emerging
classification. It is not reconstructed here. It is replaced.

\paragraph{The rule.} The universe is every area the source reports for
the indicator that is a sovereign state or dependent territory ---
operationally, an ISO~3166-1 alpha-2 code, together with Kosovo, which is
an exceptional reservation rather than a formally assigned code and is
admitted by name. Codes that are regional or income-group aggregates
rather than reporting jurisdictions are excluded and named individually in
the panel output. The rule is implemented once, in
\texttt{panel\_universe.py}, and imported by every script that needs a
country universe, so the two panel constructions cannot drift apart.

\paragraph{No history screen at this stage.} Whether an area has enough
observations to support an estimate is a property of the estimator, not of
the universe, and it is applied inside the estimation loop and reported as
attrition alongside every other reason an area drops out. Folding a
feasibility threshold into universe construction is what made the previous
list unauditable: it merged the question of which areas exist for this
purpose with the question of which areas carry enough data, and left only
the intersection visible. Keeping them separate means every exclusion is
attributable to a stated cause.

\section{Currency-union pseudo-replication: a checked and rejected
exclusion}

When the country universe was rebuilt (Section~\ref{sec:universe}), two
clusters of areas were noticed to return near-identical numbers of
observations: the Eastern Caribbean-adjacent group (Anguilla, Antigua and
Barbuda, Dominica, Grenada, St Kitts and Nevis, St Lucia, Montserrat, St
Vincent and the Grenadines) and the CEMAC-adjacent group (Cameroon,
Central African Republic, Republic of Congo, Gabon, Equatorial Guinea,
Chad). Both are currency unions with shared supervisory arrangements, so
the natural suspicion was that their members report a single aggregated
banking system, in which case treating them as independent observations
would insert eight or six near-duplicate points into any cross-country
median or dispersion statistic.

\paragraph{Decision, taken before testing.} The suspicion was recorded as
a flag in the panel output, not acted on as an exclusion. A pattern in
observation counts is not evidence about the underlying series, and
excluding countries on a plausible-sounding conjecture is precisely the
practice this document exists to prevent.

\paragraph{Test.} \texttt{panel\_currency\_union\_check.py} compares the
actual reported value series, pair by pair within each cluster: the
correlation across jointly observed quarters, and the fraction of
quarters where two members report numerically identical values. Shared
reporting would produce correlations near one together with a high
identical-value fraction.

\paragraph{Result: the suspicion is wrong.} Across all $28$ Eastern
Caribbean pairs and all $15$ CEMAC pairs, the fraction of numerically
identical quarters is \emph{zero}. Correlations are scattered across
roughly $[-0.77, +0.74]$, many of them negative, with no pair approaching
the near-unity that duplicated reporting would produce. The single
highest, Equatorial Guinea and Chad at $+0.74$ with no identical values,
is unremarkable co-movement between two small oil-exporting economies
sharing a currency and a commodity cycle --- an economic correlation, not
a reporting artefact.

\paragraph{Consequence.} These are independent national series whose
sample windows happen to begin together, which is what a common
supervisory reporting start date produces. No exclusion is warranted,
the flag has been removed from the panel script, and this section is the
record of why. The test should be re-run if the source's coverage
changes.

\section{The Basel III attempt: a documented negative result}
\label{sec:basel-negative}

Applying ground~4 to real data did not go as the diagnosis leading to it
suggested it would, and the failure is instructive enough to record in
full rather than pass over.

\paragraph{The pattern that motivated the attempt.} Constructing $z$ from
each country's own expanding-mean reference (Section~\ref{sec:units} and
its validation) and applying it to the real 77-country panel produced only
26 non-degenerate country estimates out of 77 --- a two-thirds attrition
rate against the classification guard. Among the thirteen countries with a
\emph{complete} surplus skew (zero deficit quarters observed at all), every
single one had an FSI sample window spanning the original Basel~III capital
conservation buffer phase-in, sourced to the Bank for International
Settlements' own international timetable (1 January 2013 for the CET1
step-up, phased to full effect 1 January 2019) and confirmed directly for
two of the thirteen: South Africa, via the South African Reserve Bank's own
Banks Act Circular 4 of 2016, and Indonesia, via OJK Regulation
No.\,11/POJK.03/2016. Eleven of the thirteen are BCBS, EU, or G20 members,
jurisdictions with essentially no discretion over whether to follow the
international schedule.

\paragraph{The mechanism, confirmed by simulation before touching real data
again.} A one-time, Basel-style level shift injected into an otherwise
clean simulated capital-buffer process reproduces the symptom: degeneracy
under a continuous expanding-mean reference rises from near zero to 40\%.
Segmenting the reference at the shift --- restarting the expanding mean
after the break rather than letting one continuous average straddle a
change in what ``normal capital'' means --- cuts that to 22\%. The
mechanism is real.

\paragraph{The result on real data.} Segmenting the real panel's
expanding-mean references at 1 January 2019 did not improve the usable
count: 25 non-degenerate countries against 26--27 under the unsegmented
construction, a wash within noise. More importantly, \emph{none of the
thirteen targeted countries were rescued into non-degenerate status}. Two
came close --- Latvia and the Netherlands were four and three observations
short of the minimum-regime threshold respectively --- but most were not:
Australia showed zero deficit quarters even measured against a reference
that restarts after the transition, over roughly nine years of subsequent
data. Segmentation instead lost five previously usable countries (Sri
Lanka, Macau, Nicaragua, Eswatini, T\"urkiye) while gaining three different
ones, one of which --- the United States, with an anomalously large
$\hat\lambda = 1.61$ --- is better explained as noise from a thin
post-break window than as signal.

\paragraph{Why this happened, and what it establishes.} Segmenting a
reference trades access for precision: splitting one long series into two
shorter ones repairs the classification (the split becomes balanced) but
shortens the sample each regime estimate draws on, and the resulting
small-sample attenuation offsets the gain almost exactly. Simulation
confirmed this in advance --- mean $\hat\lambda$ under segmentation was
statistically indistinguishable from the unsegmented value, both well below
the true parameter used to generate the data --- so the real-data outcome
is not a surprise in hindsight, though it was not obvious before the
attempt was made. A single calendar break, however well sourced, cannot
simultaneously fit every country's true contamination structure: it
resolves the mechanism for some observations and introduces a new mismatch
for others, netting out to no improvement in aggregate even where the
underlying diagnosis is correct.

\paragraph{Disposition.} Basel~III timing is not abandoned as an
explanation --- it remains the best-supported account of why a large block
of otherwise well-constructed observations classify as degenerate --- but
time-series segmentation is not the right tool to act on it. The
recommended path, taken up after the remaining exclusion grounds in this
section are applied, is to control for the transition at the
\emph{regression} stage: a post-2019 indicator or interaction term on the
cross-sectional estimates from the \emph{unsegmented} construction, which
does not cost the sample length that segmentation does. Whether that
control absorbs the pattern better than segmentation did is not yet known
and should not be assumed.

\end{document}
